\subsection{DDL -- Data Definition Language}
Die DDL umfasst alle Befehle, die die \emph{Struktur} einer Datenbank definieren und verändern. Es geht nicht um die Daten selbst, sondern um das Schema. Fast jedes Datenbankobjekt kann mit \icode{CREATE} erzeugt, mit \icode{ALTER} verändert, mit \icode{RENAME}/\icode{CHANGE} umbenannt und mit \icode{DROP} gelöscht werden.

\subsubsection{Datenbank und Tabelle anlegen}
Mit \icode{IF NOT EXISTS} wird eine Exception abgefangen, falls die Datenbank bereits existiert.

\begin{syntaxbox}[SQL]
CREATE DATABASE IF NOT EXISTS kurs_db;

CREATE TABLE Personal (
    PersonalNummer  INT UNSIGNED,      -- PK: NOT NULL implizit
    Vorname         VARCHAR(60)  NOT NULL,
    Nachname        VARCHAR(60)  NOT NULL,
    GeburtsDatum    DATE         NOT NULL,
    AbteilungsID    SMALLINT UNSIGNED NOT NULL,
    PRIMARY KEY (PersonalNummer),
    FOREIGN KEY (AbteilungsID) REFERENCES Abteilung(AbteilungsID)
);
\end{syntaxbox}

\subsubsection{Datentypen}
\begin{stdtable}{l X}{Art & Typen}
  Text & \icode{CHAR(n)} feste Länge \(\cdot\) \icode{VARCHAR(n)} variable Länge, maximal $n$ Zeichen \(\cdot\) \icode{TEXT}/\icode{BLOB} sehr große Inhalte \\
  Ganze Zahlen & \icode{INT}/\icode{INTEGER}, \icode{SMALLINT}, \icode{TINYINT} \\
  Gleitkommazahlen & \icode{FLOAT}, \icode{DOUBLE} -- z.\,B. \icode{FLOAT(8,2)} mit 2 Nachkommastellen \\
  Logisch & \icode{BOOL} (intern ein \icode{TINYINT}) \\
  Zeiten & \icode{DATE} \(\cdot\) \icode{TIME} \(\cdot\) \icode{DATETIME} \(\cdot\) \icode{TIMESTAMP} (setzt bei \icode{INSERT}/\icode{UPDATE} automatisch Datum und Uhrzeit) \(\cdot\) \icode{YEAR} \\
\end{stdtable}

\subsubsection{Spaltenattribute und Constraints}
\begin{stdtable}{l X}{Schlüsselwort & Bedeutung}
  \icode{PRIMARY KEY} & Primärschlüssel: eindeutig, nie \icode{NULL}, identifiziert jeden Datensatz. Kann auch aus mehreren Spalten bestehen. \\
  \icode{FOREIGN KEY} & Fremdschlüssel: verweist auf den Primärschlüssel einer anderen Tabelle. \\
  \icode{NOT NULL} & Die Spalte darf keinen \icode{NULL}-Wert enthalten. \\
  \icode{UNIQUE} & Alle Werte der Spalte müssen eindeutig sein. \\
  \icode{DEFAULT} & Standardwert, falls beim Einfügen kein Wert angegeben wird. \\
  \icode{AUTO\_INCREMENT} & Nummeriert Datensätze automatisch aufsteigend ab 1; nur für ganzzahlige, eindeutige Felder (\icode{PRIMARY KEY} oder \icode{UNIQUE}). \\
  \icode{UNSIGNED} & Nur positive Zahlen erlaubt (kein Vorzeichen). \\
\end{stdtable}

\begin{warnbox}[Stolperfalle]
In SQL kann man eine Tabelle auch ohne \icode{PRIMARY KEY} erstellen -- das sollte man aber nicht tun und den Schlüssel bereits bei der Modellierung festlegen.
\end{warnbox}

\subsubsection{Struktur ändern mit ALTER TABLE}
\icode{ALTER TABLE} verändert die Struktur einer bestehenden Tabelle, ohne die Daten zu löschen.

\begin{syntaxbox}[SQL]
-- Spalte hinzufuegen
ALTER TABLE Personal ADD Gender CHAR AFTER GeburtsDatum;
-- Datentyp aendern
ALTER TABLE Personal MODIFY GeburtsDatum DATETIME;
-- Spalte umbenennen
ALTER TABLE Personal CHANGE GeburtsDatum Geburtsdatum DATE;
-- Spalte loeschen
ALTER TABLE Personal DROP Gender;
-- Tabelle umbenennen
ALTER TABLE Personal RENAME TO Mitarbeiter;

-- Foreign Key nachtraeglich setzen
ALTER TABLE Personal
    ADD FOREIGN KEY (AbteilungsID) REFERENCES Abteilung(AbteilungsID);
\end{syntaxbox}

\subsubsection{Referenzielle Integrität}
Die Tabelle mit dem Primärschlüssel ist der \emph{Parent} (Master), die Tabelle mit dem Fremdschlüssel das \emph{Child}. Die referenzielle Integrität verhindert verwaiste Daten -- etwa eine Bestellung ohne existierenden Kunden. Mit \icode{ON DELETE} und \icode{ON UPDATE} legt man fest, was mit den Child-Datensätzen passiert:

\begin{stdtable}{l X}{Option & Verhalten}
  \icode{RESTRICT} / \icode{NO ACTION} & Löschen bzw. Ändern des Parent-Datensatzes wird verhindert. \\
  \icode{CASCADE} & Child-Datensätze werden mitgelöscht bzw. mitgeändert. \\
  \icode{SET NULL} & Der Fremdschlüssel im Child wird auf \icode{NULL} gesetzt. \\
\end{stdtable}

\subsubsection{Löschen: DROP, TRUNCATE, DELETE}
Mit \icode{DROP} wird eine Tabelle oder Datenbank komplett entfernt, mit \icode{TRUNCATE} werden nur alle Zeilen gelöscht -- die Tabelle selbst bleibt bestehen.

\begin{syntaxbox}[SQL]
DROP DATABASE kurs_db;          -- Datenbank komplett entfernen
DROP TABLE Personal;            -- Tabelle komplett entfernen
TRUNCATE TABLE Personal;        -- alle Zeilen loeschen, Struktur bleibt
\end{syntaxbox}

\begin{stdtable}{l l X}{Befehl & Löscht Daten? & Löscht Struktur?}
  \icode{DROP TABLE} & Ja & Ja -- Tabelle ist weg \\
  \icode{TRUNCATE TABLE} & Ja & Nein -- Tabelle bleibt \\
  \icode{DELETE FROM} & Ja (mit \icode{WHERE} filterbar) & Nein -- Tabelle bleibt \\
\end{stdtable}
